% Options for packages loaded elsewhere
\PassOptionsToPackage{unicode}{hyperref}
\PassOptionsToPackage{hyphens}{url}
\documentclass[
]{article}
\usepackage{xcolor}
\usepackage[margin=1in]{geometry}
\usepackage{amsmath,amssymb}
\setcounter{secnumdepth}{5}
\usepackage{iftex}
\ifPDFTeX
  \usepackage[T1]{fontenc}
  \usepackage[utf8]{inputenc}
  \usepackage{textcomp} % provide euro and other symbols
\else % if luatex or xetex
  \usepackage{unicode-math} % this also loads fontspec
  \defaultfontfeatures{Scale=MatchLowercase}
  \defaultfontfeatures[\rmfamily]{Ligatures=TeX,Scale=1}
\fi
\usepackage{lmodern}
\ifPDFTeX\else
  % xetex/luatex font selection
\fi
% Use upquote if available, for straight quotes in verbatim environments
\IfFileExists{upquote.sty}{\usepackage{upquote}}{}
\IfFileExists{microtype.sty}{% use microtype if available
  \usepackage[]{microtype}
  \UseMicrotypeSet[protrusion]{basicmath} % disable protrusion for tt fonts
}{}
\makeatletter
\@ifundefined{KOMAClassName}{% if non-KOMA class
  \IfFileExists{parskip.sty}{%
    \usepackage{parskip}
  }{% else
    \setlength{\parindent}{0pt}
    \setlength{\parskip}{6pt plus 2pt minus 1pt}}
}{% if KOMA class
  \KOMAoptions{parskip=half}}
\makeatother
\usepackage{longtable,booktabs,array}
\newcounter{none} % for unnumbered tables
\usepackage{calc} % for calculating minipage widths
% Correct order of tables after \paragraph or \subparagraph
\usepackage{etoolbox}
\makeatletter
\patchcmd\longtable{\par}{\if@noskipsec\mbox{}\fi\par}{}{}
\makeatother
% Allow footnotes in longtable head/foot
\IfFileExists{footnotehyper.sty}{\usepackage{footnotehyper}}{\usepackage{footnote}}
\makesavenoteenv{longtable}
\usepackage{graphicx}
\makeatletter
\newsavebox\pandoc@box
\newcommand*\pandocbounded[1]{% scales image to fit in text height/width
  \sbox\pandoc@box{#1}%
  \Gscale@div\@tempa{\textheight}{\dimexpr\ht\pandoc@box+\dp\pandoc@box\relax}%
  \Gscale@div\@tempb{\linewidth}{\wd\pandoc@box}%
  \ifdim\@tempb\p@<\@tempa\p@\let\@tempa\@tempb\fi% select the smaller of both
  \ifdim\@tempa\p@<\p@\scalebox{\@tempa}{\usebox\pandoc@box}%
  \else\usebox{\pandoc@box}%
  \fi%
}
% Set default figure placement to htbp
\def\fps@figure{htbp}
\makeatother
% definitions for citeproc citations
\NewDocumentCommand\citeproctext{}{}
\NewDocumentCommand\citeproc{mm}{%
  \begingroup\def\citeproctext{#2}\cite{#1}\endgroup}
\makeatletter
 % allow citations to break across lines
 \let\@cite@ofmt\@firstofone
 % avoid brackets around text for \cite:
 \def\@biblabel#1{}
 \def\@cite#1#2{{#1\if@tempswa , #2\fi}}
\makeatother
\newlength{\cslhangindent}
\setlength{\cslhangindent}{1.5em}
\newlength{\csllabelwidth}
\setlength{\csllabelwidth}{3em}
\newenvironment{CSLReferences}[2] % #1 hanging-indent, #2 entry-spacing
 {\begin{list}{}{%
  \setlength{\itemindent}{0pt}
  \setlength{\leftmargin}{0pt}
  \setlength{\parsep}{0pt}
  % turn on hanging indent if param 1 is 1
  \ifodd #1
   \setlength{\leftmargin}{\cslhangindent}
   \setlength{\itemindent}{-1\cslhangindent}
  \fi
  % set entry spacing
  \setlength{\itemsep}{#2\baselineskip}}}
 {\end{list}}
\usepackage{calc}
\newcommand{\CSLBlock}[1]{\hfill\break\parbox[t]{\linewidth}{\strut\ignorespaces#1\strut}}
\newcommand{\CSLLeftMargin}[1]{\parbox[t]{\csllabelwidth}{\strut#1\strut}}
\newcommand{\CSLRightInline}[1]{\parbox[t]{\linewidth - \csllabelwidth}{\strut#1\strut}}
\newcommand{\CSLIndent}[1]{\hspace{\cslhangindent}#1}
\setlength{\emergencystretch}{3em} % prevent overfull lines
\providecommand{\tightlist}{%
  \setlength{\itemsep}{0pt}\setlength{\parskip}{0pt}}
\usepackage{bookmark}
\IfFileExists{xurl.sty}{\usepackage{xurl}}{} % add URL line breaks if available
\urlstyle{same}
\hypersetup{
  pdftitle={Sensitivity of carryover-mitigation analysis strategies to biomarker data-generating process: a 3 x 3 simulation study},
  pdfauthor={pmsimstats team},
  hidelinks,
  pdfcreator={LaTeX via pandoc}}

\title{Sensitivity of carryover-mitigation analysis strategies to biomarker data-generating process: a 3 x 3 simulation study}
\author{pmsimstats team}
\date{2026-04-15}

\begin{document}
\maketitle

\section{Abstract}\label{abstract}

Carryover effects in aggregated N-of-1 and crossover trials degrade
the detectability of biomarker-treatment interactions. Analysis
strategies that acknowledge carryover---binary-treatment,
exposure-weighted, and lagged-treatment models---differ in how
they parameterise residual drug effect, and each makes implicit
assumptions about the functional form of the decay. We conduct a
simulation study crossing three data-generating carryover decay
forms (linear, exponential, Weibull) against three analysis-model
carryover specifications, under both mean-moderation (Architecture
A) and multivariate-normal differential-correlation (Architecture
B) biomarker data-generating processes. Results quantify the power
penalty of analysis-model mis-specification and identify which
mitigation strategy is most robust across the decay-form grid.

\section{Introduction}\label{introduction}

Biomarker-moderated treatment effects in within-subject designs
depend on how residual drug effect is modelled in both the
data-generating process (DGP) and the analysis model. The DGP
architecture comparison in the companion manuscript and in
\texttt{docs/02-dgp-mean-moderation-vs-mvn.tex} establishes that
carryover has qualitatively different power consequences under
the two architectures\textsuperscript{\citeproc{ref-hendrickson2020}{1}}. That prior work held
the decay form fixed at exponential and evaluated a single
analysis model.

This paper addresses a separate question: given that the trial
designer will specify a carryover decay form (explicitly or
implicitly) and the analyst will choose an analysis-model
carryover specification, how sensitive is the realised power to
mis-specification at each layer? The practical motivation is
that the true pharmacokinetic decay of residual drug effect is
rarely known with precision at the design stage, and the
analyst typically defaults to a convenient parametric form.

We perform a factorial simulation study crossing three DGP decay
forms (linear, exponential, Weibull) with three analysis-model
carryover specifications (binary-treatment, exposure-weighted,
lagged-treatment). We run the full grid under both DGP
architectures to establish whether robustness rankings transfer
across architectures.

\section{Simulation design}\label{simulation-design}

\subsection{Data-generating processes}\label{data-generating-processes}

Both DGP architectures from the companion manuscript are
retained:

\begin{itemize}
\tightlist
\item
  \textbf{Architecture A (direct mean moderation).} Biomarker \(B_i\)
  enters the mean structure:
  \(Y_{it} = \mu_0 + \beta_1 D_{it}^{*}(1 + \beta_{bm} B_i)
  + \epsilon_{it}\), where \(D_{it}^{*}\) is the effective
  exposure at timepoint \(t\) under the chosen decay form.
\item
  \textbf{Architecture B (MVN differential correlation).} Biomarker
  and biological response are jointly drawn from a multivariate
  normal with treatment-state-dependent correlation
  \(\text{Cor}(B_i, BR_{it}) = c_{bm} \cdot \phi(t_{sd})\), where
  \(\phi\) is the decay form.
\end{itemize}

In both architectures the decay function \(\phi : \mathbb{R}_+
\to [0, 1]\) maps time since drug discontinuation \(t_{sd}\) to a
residual-effect multiplier. We evaluate three forms
parameterised to match a common half-life \(t_{1/2}\):

\subsubsection{Linear decay}\label{linear-decay}

\[
\phi_{\text{lin}}(t_{sd}) =
  \max\left(0,\, 1 - \frac{t_{sd}}{2 t_{1/2}}\right)
\]

Residual effect vanishes completely at \(t_{sd} = 2 t_{1/2}\).
This form is biologically implausible for drugs with
first-order elimination kinetics but represents a defensible
pragmatic assumption when only the ``carryover has fully worn
off by time X'' information is available.

\subsubsection{Exponential decay}\label{exponential-decay}

\[
\phi_{\text{exp}}(t_{sd}) =
  \exp\left(-\lambda t_{sd}\right), \qquad
  \lambda = \frac{\ln 2}{t_{1/2}}
\]

The form used throughout the companion manuscript and currently
implemented in \texttt{implementations/tidyverse/R/functions.R}.
Corresponds to first-order elimination kinetics and is the
default assumption in classical pharmacokinetics.

\subsubsection{Weibull decay}\label{weibull-decay}

\[
\phi_{\text{weib}}(t_{sd}) =
  \exp\left(-\left(\lambda_w t_{sd}\right)^{k}\right), \qquad
  \lambda_w = \frac{(\ln 2)^{1/k}}{t_{1/2}}
\]

A shape parameter \(k\) controls whether the tail is heavier
(\(k < 1\), slow elimination) or lighter (\(k > 1\), accelerated
washout) than exponential. Exponential is recovered at \(k = 1\).
We report results at \(k = 0.7\) (slow elimination) and \(k = 1.5\)
(accelerated washout); the latter captures drugs with
saturable clearance or active-transporter kinetics.

\textbf{Implementation note.} At the time of writing the tidyverse
collection implements only the exponential form (see
\texttt{implementations/tidyverse/R/functions.R}, function
\texttt{apply\_carryover\_to\_component}). Adding linear and Weibull
variants is a prerequisite simulation step, tracked as a
prerequisite in \texttt{analysis/02-carryover-sensitivity/README.md}.

\subsection{Analysis-model carryover specifications}\label{analysis-model-carryover-specifications}

The analysis model is \texttt{nlme::lme(Sx\ \textasciitilde{}\ bm\ *\ X\ +\ t,\ random\ =\ \textasciitilde{}1\ \textbar{}\ ptID,\ correlation\ =\ corCAR1(form\ =\ \textasciitilde{}\ t\ \textbar{}\ ptID))} for all
specifications. The three specifications differ in the drug
exposure predictor \(X\):

\subsubsection{A1: binary-treatment (carryover ignored)}\label{a1-binary-treatment-carryover-ignored}

\(X = D_{it}\), the on-drug indicator. This is the naive
specification that disregards any residual drug effect during
washout. It serves as a reference for the damage done by
ignoring carryover.

\subsubsection{A2: exposure-weighted (matched-decay continuous predictor)}\label{a2-exposure-weighted-matched-decay-continuous-predictor}

\(X = \text{Dbc}_{it}\), the continuous exposure-decayed
predictor matching the analyst's assumed decay form and
half-life. This is the current pmsimstats-ng default
specification. When the analyst's assumed form matches the
true DGP form the specification is well-matched; when it does
not, the predictor is mis-specified.

\subsubsection{A3: lagged-treatment (estimated carryover term)}\label{a3-lagged-treatment-estimated-carryover-term}

\(X = D_{it}\) augmented with a lagged-treatment indicator
\(L_{it} = D_{i,t-1}(1 - D_{it})\) that marks off-drug
timepoints following an on-drug timepoint. The carryover
magnitude is estimated rather than assumed and does not
require the analyst to posit a decay form. Follows the
classical crossover-trial approach surveyed in\textsuperscript{\citeproc{ref-jones2014crossover}{2}}.

\subsection{Factorial grid}\label{factorial-grid}

The principal comparison is a \(3 \times 3\) factorial of DGP
decay form against analysis-model specification, replicated
under each DGP architecture:

\begin{table}

\caption{\label{tab:grid-table}Principal $3 \times 3$ simulation grid. Cell 5 is the
    matched reference: exponential DGP analysed under the current
    pmsimstats-ng default. Off-diagonal cells quantify the cost of
    DGP-analysis mis-specification. The full grid is crossed with
    Architecture $\in \{A, B\}$.}
\centering
\begin{tabular}[t]{l|l|l|l}
\hline
DGP decay & A1 binary & A2 exposure-weighted & A3 lagged-treatment\\
\hline
Linear & cell 1 & cell 2 & cell 3\\
\hline
Exponential & cell 4 & cell 5 (matched) & cell 6\\
\hline
Weibull & cell 7 & cell 8 & cell 9\\
\hline
\end{tabular}
\end{table}

\subsection{Parameter settings}\label{parameter-settings}

{\def\LTcaptype{none} % do not increment counter
\begin{longtable}[]{@{}ll@{}}
\toprule\noalign{}
Parameter & Values \\
\midrule\noalign{}
\endhead
\bottomrule\noalign{}
\endlastfoot
Biomarker moderation (\(c_{bm}\) / \(\beta_{bm}\)) & 0.0, 0.30, 0.45 \\
Carryover half-life (\(t_{1/2}\)) & 0.0, 0.5, 1.0 weeks \\
Weibull shape (\(k\)) & 0.7, 1.5 \\
Sample size (\(N\)) & 35, 70 \\
Trial design & CO, N-of-1 (Hybrid), OL+BDC \\
DGP architecture & A, B \\
Replicates per cell & 500 (target); 50 for development \\
\end{longtable}
}

The \(c_{bm} = 0\) condition serves as the type I error check
(target: nominal 5\%). The \(t_{1/2} = 0\) condition removes
carryover entirely and provides a reference for the best-case
power at each \((N, \text{design}, c_{bm})\) combination.

\subsection{Outcome metrics}\label{outcome-metrics}

For each cell we report:

\begin{enumerate}
\def\labelenumi{\arabic{enumi}.}
\tightlist
\item
  \textbf{Power} to reject \(H_0: \beta_{\text{bm:X}} = 0\) at
  \(\alpha = 0.05\).
\item
  \textbf{Type I error} at \(c_{bm} = 0\).
\item
  \textbf{Bias} of the interaction estimate: \(\hat{\beta} -
  \beta_{\text{true}}\) on the effect-size scale calibrated
  across architectures (see
  \texttt{docs/19-mean-moderation-implementation-notes.tex} for the
  calibration derivation).
\item
  \textbf{Monte Carlo SE} on all reported estimates.
\end{enumerate}

\subsection{Computing environment}\label{computing-environment}

Simulations are orchestrated from
\texttt{analysis/02-carryover-sensitivity/} and use
\texttt{implementations/original-extended/} (data.table, dual
architecture) for primary runs. Tidyverse runs via
\texttt{implementations/tidyverse/} serve as cross-validation.
Reproducibility metadata (R version, \texttt{renv.lock} hash,
\texttt{implementations/} commit SHA, simulation seed) are recorded in
each output \texttt{.rds} file.

Parallel execution uses \texttt{furrr::future\_map} with the \texttt{multicore}
backend. Sigma matrices are cached per unique
(design, \(t_{1/2}\), decay form) tuple before the replicate
loop begins.

\section{Results}\label{results}

\subsection{Matched-decay baseline}\label{matched-decay-baseline}

Cell 5 of Table 1 (exponential DGP, exposure-weighted analysis)
is the matched-decay reference. Power under this cell matches
the Architecture B numbers reported in the companion manuscript
for the same \((N, \text{design}, c_{bm}, t_{1/2})\) setting,
validating the simulation pipeline.

\subsection{Cost of analysis-model mis-specification}\label{cost-of-analysis-model-mis-specification}

\emph{(Pending. Off-diagonal cells quantify the power penalty of
assuming the wrong decay form in the analysis model when the
DGP is fixed. Presented as a heatmap: rows = DGP form, columns =
analysis specification, cells = relative power vs.~the matched
cell for that DGP row.)}

\subsection{Cost of DGP mis-specification under a fixed analysis}\label{cost-of-dgp-mis-specification-under-a-fixed-analysis}

\emph{(Pending. Holding the analysis specification fixed at A2
(exposure-weighted exponential, the pmsimstats-ng default), the
power loss when the true DGP is linear or Weibull relative to
exponential. This is the quantity most relevant to a trialist
reporting power under the default analysis.)}

\subsection{Robustness of the lagged-treatment specification}\label{robustness-of-the-lagged-treatment-specification}

\emph{(Pending. The A3 column of the grid quantifies whether
estimating carryover as a nuisance is more robust than
assuming a specific parametric form. Classical crossover-trial
intuition\textsuperscript{\citeproc{ref-jones2014crossover}{2}} predicts A3 is more robust to
decay-form mis-specification but less efficient when the
assumed form is correct; simulation will quantify the
trade-off.)}

\subsection{Architecture interaction}\label{architecture-interaction}

\emph{(Pending. Full grid under Architecture A versus Architecture
B. Key question: do the robustness rankings across analysis
specifications transfer between architectures, or is the best
analysis specification architecture-dependent?)}

\subsection{Type I error control}\label{type-i-error-control}

\emph{(Pending. \(c_{bm} = 0\) grid. Target: all cells nominal.
Inflation would indicate that the analysis model is
mis-specified in a way that biases the reference distribution
of the interaction test.)}

\section{Discussion}\label{discussion}

\emph{(Pending. Principal findings:}

\begin{itemize}
\tightlist
\item
  \emph{Which analysis specification is most robust across DGP decay
  forms?}
\item
  \emph{Is the architecture-dependent ranking of Section 6 of the
  companion manuscript (which was derived under a single decay
  form) stable when decay form varies?}
\item
  \emph{Recommended default for a trialist who will specify the
  analysis model prospectively without knowing the true decay
  form.}
\item
  \emph{Relation to the latent-subtype tension raised by reviewer RH
  in the companion manuscript: if the true DGP is closer to an
  explicit mixture model than to any of the three parametric
  decay forms evaluated here, the lagged-treatment
  specification's non-parametric character may be doubly
  attractive.}
\item
  \emph{Limitations: we do not evaluate trial-design-level
  mitigations (longer washout, design modification; see
  \texttt{docs/02} Section 6.6) because these are orthogonal to the
  analysis-model question addressed here.)}
\end{itemize}

\section{Conclusions}\label{conclusions}

\emph{(Pending. Reporting recommendations for trialists planning a
biomarker-moderated N-of-1 study when carryover is anticipated:
choice of decay-form assumption, choice of analysis-model
specification, and sensitivity analyses to report alongside the
primary power calculation.)}

\section*{References}\label{references}
\addcontentsline{toc}{section}{References}

\protect\phantomsection\label{refs}
\begin{CSLReferences}{0}{1}
\bibitem[\citeproctext]{ref-hendrickson2020}
\CSLLeftMargin{1. }%
\CSLRightInline{Hendrickson RC, Thomas RG, Schork NJ, Raskind MA. Optimizing aggregated {N}-of-1 trial designs for predictive biomarker validation: Statistical methods and practical considerations. \emph{Frontiers in Digital Health}. 2020;2:13.}

\bibitem[\citeproctext]{ref-jones2014crossover}
\CSLLeftMargin{2. }%
\CSLRightInline{Jones B, Kenward MG. \emph{Design and Analysis of Cross-over Trials}. 3rd ed. CRC Press; 2014.}

\end{CSLReferences}

\end{document}
